%! Setting Up a Test for a Population Mean — Unit 7, Lesson 7.4 — Group Activity
\documentclass[10pt]{article}
\usepackage{apstats-article}
\usepackage{apstats-boxes}
\newcommand{\TallMath}[1]{$\displaystyle #1\rule[-1.4em]{0pt}{3.2em}$}

\begin{document}

\pageheader{Unit 7, Lesson 7.4}{Group Activity}
\namepartnerperiod

\vspace{0.15in}

\textit{For each scenario, (a) name the appropriate inference method, (b) state $H_0$ and
$H_a$ (including the direction of $H_a$), and (c) verify both conditions. Do \textbf{not}
compute a test statistic yet.}

\begin{scenariobox}[Scenario A --- Sleep Duration]{sky}
\textbf{Tier A \& R.} \quad
The American Academy of Sleep Medicine recommends that teenagers get at least 8 hours of sleep
per night. A school counselor believes students at her school sleep \emph{less} than 8 hours
on average. She randomly selects 20 students and records their nightly sleep duration (hours).
The dotplot shows the distribution is roughly symmetric with no outliers.

\begin{enumerate}[label=\alph*.]
  \item Inference method: \blank{5cm} \quad $df = $ \blank{1.5cm}
  \item $H_0$: \blank{4.5cm} \quad $H_a$: \blank{4.5cm}

        Direction of $H_a$: \blank{2cm} \quad Justify: \writeline

  \item Conditions:
        \begin{itemize}
          \item Random: \writeline
          \item Normal ($n = 20 < 30$): \writeline
        \end{itemize}
\end{enumerate}
\end{scenariobox}

\begin{scenariobox}[Scenario B --- Package Weights]{sky}
\textbf{Tier A \& R.} \quad
A snack company labels its chip bags ``Net weight: 1.5 oz.'' A quality-control inspector
suspects the \emph{true mean} weight differs from 1.5 oz. She randomly selects 40 bags and
weighs each one. The histogram of weights appears approximately bell-shaped.

\begin{enumerate}[label=\alph*.]
  \item Inference method: \blank{5cm} \quad $df = $ \blank{1.5cm}
  \item $H_0$: \blank{4.5cm} \quad $H_a$: \blank{4.5cm}

        Direction of $H_a$: \blank{2cm} \quad Justify: \writeline

  \item Conditions:
        \begin{itemize}
          \item Random: \writeline
          \item Normal ($n = 40 \ge 30$): \writeline
        \end{itemize}
\end{enumerate}
\end{scenariobox}

\begin{scenariobox}[Scenario C --- Matched Pairs: Reading Fluency]{sky}
\textbf{Tier A \& E.} \quad
A reading specialist administers a fluency test to 12 students at the beginning and end of a
semester-long intervention. She wants to determine whether the intervention \emph{improved}
reading speed (words per minute). The differences (end $-$ begin) are available for analysis.
A boxplot of the differences shows no strong skewness or outliers.

\begin{enumerate}[label=\alph*.]
  \item What type of data is this? Circle: \quad one sample \quad matched pairs

  \item Define the parameter $\mu_d$ in context: \writeline

  \item Define the order of subtraction and state hypotheses:

        $d_i = $ \blank{3cm} $-$ \blank{3cm} \qquad $H_0$: \blank{3.5cm} \quad $H_a$: \blank{3.5cm}

  \item Verify conditions:
        \begin{itemize}
          \item Random/Independence: \writeline
          \item Normal ($n = 12 < 30$): \writeline
        \end{itemize}
\end{enumerate}
\end{scenariobox}

\begin{reflectionbox}
Scenarios A and B both involve unknown $\sigma$. What is the key difference in the Normal
condition check between them? Why does sample size matter here?
\writelines{2}
\end{reflectionbox}

\end{document}
